% EXPECTED-LEAKS: settopmatter, acmConference, copyrightyear, affiliation, institution, ccsdesc, additionalaffiliation, country, city
\documentclass[sigchi]{acmart}

\settopmatter{printacmref=false}

\copyrightyear{2026}
\acmConference[CHI '26]{ACM CHI 2026}{May 4--9, 2026}{Honolulu}

\begin{document}

\title{Ambient Voice Assistants in Domestic Settings}

\author{User Researcher}
\affiliation{%
  \institution{Example HCI Lab}
  \city{Boston}
  \country{USA}
}

\begin{abstract}
We report a diary study of 22 families using ambient voice assistants over 12 weeks. We find routine breakdowns and trust erosion driven by mishearings.
\end{abstract}

\ccsdesc[500]{Human-centered computing~Empirical studies in HCI}

\keywords{voice assistants, ambient computing, diary study}

\maketitle

\section{Introduction}
Voice assistants are now ubiquitous in domestic spaces. Their ambient, always-listening nature raises open questions about reliability and privacy.

\section{Method}
We conducted a 12-week diary study with 22 households recruited through local community boards.

\section{Findings}
Participants reported an average of 3.4 mishearings per week per household. Trust in the device dropped significantly when mishearings had consequential actions (e.g.\ unintended purchases).

\section{Discussion}
Design implications include more conservative defaults for consequential actions and better transparency about what was heard.

\section{Conclusion}
Ambient assistants require design attention to the failure path, not just the success path.

\end{document}
